\documentclass[11pt]{article}

\usepackage[a4paper,margin=1in]{geometry}
\usepackage{amsmath,amssymb,amsthm,mathtools}
\usepackage{mathrsfs}
\usepackage{hyperref}
\usepackage{enumitem}
\usepackage{microtype}
\usepackage{bm}

\hypersetup{
	colorlinks=true,
	linkcolor=blue,
	urlcolor=blue,
	citecolor=blue
}

\title{Lecture Notes: Counting \(k\)-Dimensional Subspaces of an \(n\)-Dimensional Vector Space via Linear Operators}
\author{}
\date{}

\newtheorem{theorem}{Theorem}[section]
\newtheorem{proposition}[theorem]{Proposition}
\newtheorem{lemma}[theorem]{Lemma}
\newtheorem{corollary}[theorem]{Corollary}
\theoremstyle{definition}
\newtheorem{definition}[theorem]{Definition}
\newtheorem{example}[theorem]{Example}
\newtheorem{exercise}[theorem]{Exercise}
\theoremstyle{remark}
\newtheorem{remark}[theorem]{Remark}

\newcommand{\F}{\mathbf{F}}
\newcommand{\R}{\mathbf{R}}
\newcommand{\C}{\mathbf{C}}
\newcommand{\Q}{\mathbf{Q}}
\newcommand{\GL}{\mathrm{GL}}
\newcommand{\End}{\mathrm{End}}
\newcommand{\Hom}{\mathrm{Hom}}
\newcommand{\rank}{\operatorname{rank}}
\newcommand{\nullity}{\operatorname{nullity}}
\newcommand{\im}{\operatorname{im}}
\newcommand{\Ker}{\operatorname{Ker}}
\newcommand{\Span}{\operatorname{span}}
\newcommand{\Gr}{\mathrm{Gr}}
\newcommand{\qbinom}[2]{\genfrac{[}{]}{0pt}{}{#1}{#2}_q}

\begin{document}
	
	\maketitle
	
	\tableofcontents
	
	\section{Introduction}
	
	Let \(V\) be an \(n\)-dimensional vector space over a field \(K\). We want to answer the question:
	
	\begin{quote}
		How many distinct \(k\)-dimensional subspaces does \(V\) have?
	\end{quote}
	
	If \(K\) is infinite, then for \(0<k<n\) the answer is infinite. The interesting counting problem occurs when \(K=\F_q\) is a finite field with \(q\) elements.
	
	The standard answer is:
	\[
	\#\{\text{\(k\)-dimensional subspaces of }\F_q^n\}=\qbinom{n}{k}.
	\]
	
	There are several ways to prove this. In these notes, we develop a proof and interpretation in terms of \emph{linear operators}. More precisely, we study surjective linear operators
	\[
	T:\F_q^n\to \F_q^k
	\]
	and show that the kernels of such operators encode the \((n-k)\)-dimensional subspaces of \(\F_q^n\). By counting such operators carefully, we recover the Gaussian binomial coefficient.
	
	So the philosophy is:
	
	\begin{quote}
		Subspaces can be studied through kernels, images, and ranks of linear operators.
	\end{quote}
	
	This viewpoint is very natural from linear algebra, since subspaces often arise as kernels and images of maps.
	
	\section{The Main Result}
	
	\begin{definition}
		Let \(0\le k\le n\). The \emph{Gaussian binomial coefficient} is defined by
		\[
		\qbinom{n}{k}
		=
		\frac{(q^n-1)(q^n-q)(q^n-q^2)\cdots(q^n-q^{k-1})}
		{(q^k-1)(q^k-q)(q^k-q^2)\cdots(q^k-q^{k-1})}.
		\]
		Equivalently,
		\[
		\qbinom{n}{k}
		=
		\prod_{i=0}^{k-1}\frac{q^n-q^i}{q^k-q^i}.
		\]
	\end{definition}
	
	\begin{theorem}[Main Theorem]
		Let \(V\) be an \(n\)-dimensional vector space over \(\F_q\). Then the number of distinct \(k\)-dimensional subspaces of \(V\) is
		\[
		\boxed{
			\qbinom{n}{k}
			=
			\prod_{i=0}^{k-1}\frac{q^n-q^i}{q^k-q^i}
		}.
		\]
	\end{theorem}
	
	Our goal is to prove this formula using linear operators.
	
	\section{Why Linear Operators Are Relevant}
	
	A subspace can often be described as a kernel or as an image.
	
	\begin{itemize}
		\item If \(T:V\to W\) is linear, then \(\Ker(T)\subseteq V\) is a subspace.
		\item If \(T:V\to W\) is linear, then \(\im(T)\subseteq W\) is a subspace.
	\end{itemize}
	
	Thus one can try to count subspaces by counting linear operators with prescribed kernel or image.
	
	For our counting problem, kernels are especially useful. If
	\[
	T:\F_q^n\to \F_q^k
	\]
	is surjective, then by rank-nullity,
	\[
	\dim\Ker(T)=n-k.
	\]
	So surjective linear maps from \(\F_q^n\) onto \(\F_q^k\) produce \((n-k)\)-dimensional subspaces of \(\F_q^n\).
	
	Conversely, every \((n-k)\)-dimensional subspace arises as the kernel of many surjective operators.
	
	This gives a counting strategy:
	\begin{enumerate}[label=\arabic*.]
		\item count surjective linear operators \(T:\F_q^n\to \F_q^k\),
		\item determine how many such operators have a fixed kernel,
		\item divide.
	\end{enumerate}
	
	That is the operator-theoretic version of the standard basis-counting argument.
	
	\section{Preliminaries on Linear Operators}
	
	We recall the fundamental dimension theorem.
	
	\begin{theorem}[Rank-Nullity]
		Let \(T:V\to W\) be a linear operator between finite-dimensional vector spaces. Then
		\[
		\dim V=\rank(T)+\nullity(T),
		\]
		that is,
		\[
		\dim V=\dim\im(T)+\dim\Ker(T).
		\]
	\end{theorem}
	
	\begin{proof}
		Choose a basis \(u_1,\dots,u_r\) of \(\Ker(T)\), and extend it to a basis
		\[
		u_1,\dots,u_r,v_1,\dots,v_s
		\]
		of \(V\). Then one checks that
		\[
		T(v_1),\dots,T(v_s)
		\]
		form a basis of \(\im(T)\). Hence
		\[
		\dim V=r+s=\dim\Ker(T)+\dim\im(T).
		\]
	\end{proof}
	
	\begin{corollary}
		If \(T:\F_q^n\to \F_q^k\) is surjective, then
		\[
		\dim\Ker(T)=n-k.
		\]
	\end{corollary}
	
	\begin{proof}
		If \(T\) is surjective, then \(\rank(T)=k\). By rank-nullity,
		\[
		n=\rank(T)+\nullity(T)=k+\nullity(T),
		\]
		so \(\nullity(T)=n-k\).
	\end{proof}
	
	So every surjective linear operator \(\F_q^n\to \F_q^k\) determines an \((n-k)\)-dimensional subspace, namely its kernel.
	
	\section{Every Subspace Is a Kernel}
	
	The next fact explains why kernels capture all subspaces of the appropriate dimension.
	
	\begin{proposition}
		Let \(U\subseteq \F_q^n\) be a subspace of dimension \(n-k\). Then there exists a surjective linear operator
		\[
		T:\F_q^n\to \F_q^k
		\]
		such that
		\[
		\Ker(T)=U.
		\]
	\end{proposition}
	
	\begin{proof}
		Choose a basis
		\[
		u_1,\dots,u_{n-k}
		\]
		of \(U\), and extend it to a basis
		\[
		u_1,\dots,u_{n-k},v_1,\dots,v_k
		\]
		of \(\F_q^n\). Define \(T\) on this basis by
		\[
		T(u_i)=0 \quad (1\le i\le n-k), \qquad T(v_j)=e_j \quad (1\le j\le k),
		\]
		where \(e_1,\dots,e_k\) is the standard basis of \(\F_q^k\). Extend linearly.
		
		Then \(T\) is surjective because the \(e_j\)'s lie in its image, and
		\[
		\Ker(T)=U.
		\]
	\end{proof}
	
	\begin{remark}
		Thus \((n-k)\)-dimensional subspaces of \(\F_q^n\) are exactly the kernels of surjective operators
		\[
		\F_q^n\to \F_q^k.
		\]
	\end{remark}
	
	Since counting \(k\)-dimensional subspaces is the same as counting \((n-k)\)-dimensional subspaces, kernels already solve the whole problem in principle.
	
	\section{Counting Surjective Operators}
	
	We now count the number of surjective linear operators
	\[
	T:\F_q^n\to \F_q^k.
	\]
	
	A linear operator is determined by the images of a basis of the domain. Equivalently, after choosing bases, it is represented by a \(k\times n\) matrix over \(\F_q\). The operator is surjective exactly when the matrix has rank \(k\).
	
	\begin{proposition}
		The number of surjective linear operators
		\[
		T:\F_q^n\to \F_q^k
		\]
		is
		\[
		(q^n-1)(q^n-q)(q^n-q^2)\cdots(q^n-q^{k-1}).
		\]
	\end{proposition}
	
	\begin{proof}
		A surjective map \(T:\F_q^n\to \F_q^k\) has rank \(k\). The rows of a representing \(k\times n\) matrix must therefore be linearly independent in the dual viewpoint, or equivalently the columns must span \(\F_q^k\). It is simplest to count by choosing the rows as an ordered linearly independent \(k\)-tuple in \((\F_q^n)^*\), which has the same count as an ordered linearly independent \(k\)-tuple in \(\F_q^n\).
		
		Thus:
		\begin{itemize}
			\item the first row can be any nonzero vector: \(q^n-1\) choices,
			\item the second row must be outside its span: \(q^n-q\) choices,
			\item the third row must be outside the span of the first two: \(q^n-q^2\) choices,
		\end{itemize}
		and so on.
		
		Hence the number of surjective operators is
		\[
		(q^n-1)(q^n-q)(q^n-q^2)\cdots(q^n-q^{k-1}).
		\]
	\end{proof}
	
	\begin{remark}
		This is also the number of \(k\times n\) matrices of rank \(k\).
	\end{remark}
	
	\section{How Many Surjective Operators Have a Fixed Kernel?}
	
	Now fix an \((n-k)\)-dimensional subspace
	\[
	U\subseteq \F_q^n.
	\]
	We ask: how many surjective linear operators
	\[
	T:\F_q^n\to \F_q^k
	\]
	satisfy
	\[
	\Ker(T)=U?
	\]
	
	The answer is controlled by the quotient space \(\F_q^n/U\).
	
	\begin{proposition}
		Let \(U\subseteq \F_q^n\) be a subspace of dimension \(n-k\). Then the number of surjective linear operators
		\[
		T:\F_q^n\to \F_q^k
		\]
		with kernel \(U\) is
		\[
		|\GL_k(\F_q)|
		=
		(q^k-1)(q^k-q)(q^k-q^2)\cdots(q^k-q^{k-1}).
		\]
	\end{proposition}
	
	\begin{proof}
		Since \(\dim U=n-k\), the quotient space \(\F_q^n/U\) has dimension
		\[
		n-(n-k)=k.
		\]
		If \(T:\F_q^n\to \F_q^k\) has kernel \(U\), then by the universal property of quotients, \(T\) factors uniquely as
		\[
		\F_q^n \xrightarrow{\pi} \F_q^n/U \xrightarrow{\widetilde{T}} \F_q^k,
		\]
		where \(\pi\) is the quotient map and \(\widetilde{T}\) is linear.
		
		Now \(T\) is surjective and \(\dim(\F_q^n/U)=\dim(\F_q^k)=k\), so \(\widetilde{T}\) must be an isomorphism. Conversely, every isomorphism
		\[
		\widetilde{T}:\F_q^n/U\to \F_q^k
		\]
		gives a surjective linear map \(T=\widetilde{T}\circ\pi\) with kernel \(U\).
		
		So the number of such \(T\) is exactly the number of vector space isomorphisms
		\[
		\F_q^n/U\to \F_q^k,
		\]
		which is the number of ordered bases of \(\F_q^k\), namely
		\[
		|\GL_k(\F_q)|
		=
		(q^k-1)(q^k-q)(q^k-q^2)\cdots(q^k-q^{k-1}).
		\]
	\end{proof}
	
	\begin{remark}
		This is the central linear-operator idea: a fixed kernel corresponds not to one operator, but to all possible isomorphisms from the quotient \(V/U\) onto the codomain.
	\end{remark}
	
	\section{The Counting Formula from Kernels of Operators}
	
	We are now ready for the main counting argument.
	
	\begin{proof}[Proof of the Main Theorem]
		Let
		\[
		\mathcal{S}=\{T:\F_q^n\to \F_q^k \mid T \text{ is surjective}\}.
		\]
		We count \(|\mathcal{S}|\) in two ways.
		
		First, by the proposition above,
		\[
		|\mathcal{S}|=(q^n-1)(q^n-q)(q^n-q^2)\cdots(q^n-q^{k-1}).
		\]
		
		Second, each surjective operator \(T\) has kernel of dimension \(n-k\), and every \((n-k)\)-dimensional subspace arises as such a kernel. For each fixed \((n-k)\)-dimensional subspace \(U\), the number of surjective operators with kernel \(U\) is
		\[
		|\GL_k(\F_q)|
		=
		(q^k-1)(q^k-q)(q^k-q^2)\cdots(q^k-q^{k-1}).
		\]
		
		Therefore,
		\[
		|\mathcal{S}|
		=
		\bigl(\#\{\text{\((n-k)\)-dimensional subspaces of }\F_q^n\}\bigr)\cdot |\GL_k(\F_q)|.
		\]
		
		So
		\[
		\#\{\text{\((n-k)\)-dimensional subspaces of }\F_q^n\}
		=
		\frac{(q^n-1)(q^n-q)\cdots(q^n-q^{k-1})}
		{(q^k-1)(q^k-q)\cdots(q^k-q^{k-1})}.
		\]
		
		By symmetry of Gaussian coefficients,
		\[
		\qbinom{n}{n-k}=\qbinom{n}{k},
		\]
		so this is also the number of \(k\)-dimensional subspaces. Hence
		\[
		\#\{\text{\(k\)-dimensional subspaces of }\F_q^n\}
		=
		\qbinom{n}{k}.
		\]
	\end{proof}
	
	\section{A Direct Operator Form of the Answer}
	
	The previous proof can be summarized in a clean operator identity:
	
	\begin{align*}
		\#\{\text{\(k\)-dimensional subspaces of }\F_q^n\}
		&=
		\#\{\text{\((n-k)\)-dimensional subspaces of }\F_q^n\} \\
		&=
		\frac{\#\{\text{surjective }T:\F_q^n\to \F_q^k\}}
		{\#\{\text{automorphisms of }\F_q^k\}} \\
		&=
		\frac{\#\{\text{surjective }T:\F_q^n\to \F_q^k\}}
		{|\GL_k(\F_q)|}.
	\end{align*}
	
	This is the operator-theoretic analogue of dividing by the number of bases of a fixed subspace.
	
	\section{Images Instead of Kernels}
	
	One may also formulate the same story using images.
	
	Let
	\[
	S:\F_q^k\to \F_q^n
	\]
	be an injective linear operator. Then \(\im(S)\) is a \(k\)-dimensional subspace of \(\F_q^n\). Conversely, every \(k\)-dimensional subspace arises as the image of many injective operators.
	
	\begin{proposition}
		The number of injective linear operators
		\[
		S:\F_q^k\to \F_q^n
		\]
		is
		\[
		(q^n-1)(q^n-q)(q^n-q^2)\cdots(q^n-q^{k-1}).
		\]
	\end{proposition}
	
	\begin{proof}
		To define an injective map \(S\), choose the images of a basis of \(\F_q^k\). These images must be linearly independent in \(\F_q^n\). So we count ordered linearly independent \(k\)-tuples in \(\F_q^n\), obtaining
		\[
		(q^n-1)(q^n-q)\cdots(q^n-q^{k-1}).
		\]
	\end{proof}
	
	\begin{proposition}
		If \(W\subseteq \F_q^n\) is a fixed \(k\)-dimensional subspace, then the number of injective operators
		\[
		S:\F_q^k\to \F_q^n
		\]
		with image \(W\) is
		\[
		|\GL_k(\F_q)|.
		\]
	\end{proposition}
	
	\begin{proof}
		Such an injective operator is exactly an isomorphism
		\[
		\F_q^k \xrightarrow{\sim} W.
		\]
		The number of these is the number of ordered bases of \(W\), namely \(|\GL_k(\F_q)|\).
	\end{proof}
	
	Therefore,
	\[
	\#\{\text{\(k\)-dimensional subspaces of }\F_q^n\}
	=
	\frac{\#\{\text{injective }S:\F_q^k\to \F_q^n\}}
	{|\GL_k(\F_q)|}.
	\]
	
	So kernels and images give two dual operator proofs of the same formula.
	
	\section{Change of Basis and Operator Invariance}
	
	It is important that the answer does not depend on a chosen basis.
	
	\begin{proposition}
		If \(V\) and \(W\) are isomorphic finite-dimensional vector spaces over \(\F_q\), then the number of \(k\)-dimensional subspaces of \(V\) equals the number of \(k\)-dimensional subspaces of \(W\).
	\end{proposition}
	
	\begin{proof}
		If \(A:V\to W\) is an isomorphism, then
		\[
		U\longmapsto A(U)
		\]
		is a bijection between the \(k\)-dimensional subspaces of \(V\) and the \(k\)-dimensional subspaces of \(W\).
	\end{proof}
	
	Thus we may always choose a basis and work in \(\F_q^n\). The operator counts above are intrinsic, because rank, kernel, image, and isomorphism type are basis-independent notions.
	
	\section{Examples}
	
	\begin{example}[Lines in \(\F_q^3\)]
		The number of \(1\)-dimensional subspaces is
		\[
		\qbinom{3}{1}
		=
		\frac{q^3-1}{q-1}
		=
		q^2+q+1.
		\]
		Operator interpretation: these are precisely the images of injective maps
		\[
		\F_q\to \F_q^3,
		\]
		modulo automorphisms of \(\F_q\), or equivalently the kernels of surjective maps
		\[
		\F_q^3\to \F_q^2.
		\]
	\end{example}
	
	\begin{example}[Planes in \(\F_q^3\)]
		The number of \(2\)-dimensional subspaces is
		\[
		\qbinom{3}{2}=\qbinom{3}{1}=q^2+q+1.
		\]
		Operator interpretation: these are the kernels of nonzero linear functionals
		\[
		T:\F_q^3\to \F_q.
		\]
		Each nonzero functional is surjective, and two such functionals have the same kernel exactly when they differ by a nonzero scalar. Hence
		\[
		\frac{q^3-1}{q-1}=q^2+q+1.
		\]
	\end{example}
	
	\begin{example}[A very concrete operator count for \(\F_2^3\)]
		Let us count the \(2\)-dimensional subspaces of \(\F_2^3\).
		
		A \(2\)-dimensional subspace is the kernel of a surjective operator
		\[
		T:\F_2^3\to \F_2.
		\]
		Since \(\F_2\) has only one nonzero scalar, every nonzero linear functional has a distinct kernel count contribution of size \(1\). There are
		\[
		2^3-1=7
		\]
		nonzero linear functionals, hence exactly \(7\) distinct kernels, i.e. \(7\) planes in \(\F_2^3\).
		
		This agrees with
		\[
		\qbinom{3}{2}\Big|_{q=2}=7.
		\]
	\end{example}
	
	\section{Special Cases}
	
	\begin{proposition}
		For every \(n\) and \(q\),
		\begin{align*}
			\qbinom{n}{0} &= 1,\\
			\qbinom{n}{1} &= \frac{q^n-1}{q-1}=1+q+q^2+\cdots+q^{n-1},\\
			\qbinom{n}{n-1} &= \qbinom{n}{1},\\
			\qbinom{n}{n} &= 1.
		\end{align*}
	\end{proposition}
	
	\begin{proof}
		The cases \(k=0\) and \(k=n\) are obvious. For \(k=1\), count injective maps
		\[
		\F_q\to \F_q^n
		\]
		or equivalently nonzero vectors in \(\F_q^n\), and divide by the nonzero scalars \(q-1\). This gives
		\[
		\qbinom{n}{1}=\frac{q^n-1}{q-1}.
		\]
		The case \(k=n-1\) follows dually by kernels of nonzero linear functionals.
	\end{proof}
	
	\section{Symmetry}
	
	\begin{proposition}
		For \(0\le k\le n\),
		\[
		\qbinom{n}{k}=\qbinom{n}{n-k}.
		\]
	\end{proposition}
	
	\begin{proof}
		One proof is algebraic from the product formula. Another proof comes from operators: \(k\)-dimensional subspaces can be counted by images of injective maps
		\[
		\F_q^k\to \F_q^n,
		\]
		while \((n-k)\)-dimensional subspaces can be counted by kernels of surjective maps
		\[
		\F_q^n\to \F_q^k.
		\]
		These yield the same quotient formula.
	\end{proof}
	
	\section{Relation with Rank-\(k\) Operators}
	
	There is a broader operator viewpoint.
	
	A linear operator
	\[
	T:\F_q^n\to \F_q^n
	\]
	of rank \(k\) has:
	\begin{itemize}
		\item image a \(k\)-dimensional subspace,
		\item kernel an \((n-k)\)-dimensional subspace.
	\end{itemize}
	
	So rank-\(k\) operators simultaneously encode two subspaces. This connects the counting of subspaces with the classification of linear maps by rank.
	
	One can therefore think of Gaussian coefficients as arising naturally from the structure theory of finite-dimensional operators over finite fields.
	
	\section{Grassmannians from the Operator Viewpoint}
	
	The set of all \(k\)-dimensional subspaces of \(\F_q^n\) is called the Grassmannian:
	\[
	\Gr(k,n).
	\]
	
	From the operator viewpoint, \(\Gr(k,n)\) can be realized in two ways:
	
	\begin{itemize}
		\item as the set of images of injective operators
		\[
		\F_q^k \to \F_q^n
		\]
		modulo automorphisms of \(\F_q^k\),
		\item as the set of kernels of surjective operators
		\[
		\F_q^n \to \F_q^{n-k}
		\]
		modulo automorphisms of the codomain.
	\end{itemize}
	
	This makes the Grassmannian a natural moduli space for linear operators of full rank.
	
	\section{Infinite Fields}
	
	\begin{proposition}
		If \(K\) is infinite and \(V\) is an \(n\)-dimensional vector space over \(K\), then for \(0<k<n\) there are infinitely many \(k\)-dimensional subspaces of \(V\).
	\end{proposition}
	
	\begin{proof}
		For \(k=1\), every line through the origin in \(K^2\) is a \(1\)-dimensional subspace, and there are infinitely many such lines when \(K\) is infinite. Hence there are infinitely many \(k\)-dimensional subspaces in general for \(0<k<n\).
	\end{proof}
	
	\begin{remark}
		So the finite formula
		\[
		\qbinom{n}{k}
		\]
		is special to finite fields.
	\end{remark}
	
	\section{Conceptual Summary}
	
	The operator-theoretic proof rests on the following facts:
	
	\begin{enumerate}[label=\arabic*.]
		\item Every \((n-k)\)-dimensional subspace \(U\subseteq \F_q^n\) is the kernel of a surjective operator
		\[
		T:\F_q^n\to \F_q^k.
		\]
		\item The number of surjective operators \(\F_q^n\to \F_q^k\) is
		\[
		(q^n-1)(q^n-q)\cdots(q^n-q^{k-1}).
		\]
		\item For a fixed kernel \(U\), the number of such operators is
		\[
		|\GL_k(\F_q)|
		=
		(q^k-1)(q^k-q)\cdots(q^k-q^{k-1}).
		\]
		\item Dividing gives the number of subspaces.
	\end{enumerate}
	
	So:
	\[
	\boxed{
		\#\{\text{\(k\)-dimensional subspaces of }\F_q^n\}
		=
		\frac{\#\{\text{surjective }T:\F_q^n\to \F_q^k\}}
		{|\GL_k(\F_q)|}
		=
		\qbinom{n}{k}
	}.
	\]
	
	\section{Final Theorem}
	
	\begin{theorem}[Final Form]
		Let \(V\) be an \(n\)-dimensional vector space over \(\F_q\). Then the number of distinct \(k\)-dimensional subspaces of \(V\) is
		\[
		\boxed{
			\qbinom{n}{k}
			=
			\prod_{i=0}^{k-1}\frac{q^n-q^i}{q^k-q^i}
		}.
		\]
		Equivalently, this number can be computed from linear operators by
		\[
		\#\{\text{\(k\)-dimensional subspaces of }V\}
		=
		\frac{\#\{\text{injective }S:\F_q^k\to V\}}{|\GL_k(\F_q)|},
		\]
		or dually by
		\[
		\#\{\text{\(k\)-dimensional subspaces of }V\}
		=
		\frac{\#\{\text{surjective }T:V\to \F_q^k\}}{|\GL_k(\F_q)|}.
		\]
	\end{theorem}
	
	\section{Exercises}
	
	\begin{exercise}
		Let \(U\subseteq \F_q^n\) have dimension \(n-k\). Construct explicitly a surjective linear operator
		\[
		T:\F_q^n\to \F_q^k
		\]
		with \(\Ker(T)=U\).
	\end{exercise}
	
	\begin{exercise}
		Prove carefully that if \(U\subseteq \F_q^n\) has dimension \(n-k\), then the number of surjective linear operators \(T:\F_q^n\to \F_q^k\) with kernel \(U\) is \(|\GL_k(\F_q)|\).
	\end{exercise}
	
	\begin{exercise}
		Count the number of injective maps
		\[
		\F_q^2\to \F_q^4.
		\]
		Use this to compute \(\qbinom{4}{2}\).
	\end{exercise}
	
	\begin{exercise}
		Show directly that the number of nonzero linear functionals
		\[
		T:\F_q^n\to \F_q
		\]
		modulo nonzero scalar multiples is \(\qbinom{n}{n-1}\).
	\end{exercise}
	
	\begin{exercise}
		Compute the number of \(2\)-dimensional subspaces of \(\F_2^4\).
	\end{exercise}
	
	\begin{exercise}
		Explain why images of injective maps and kernels of surjective maps give the same counting formula.
	\end{exercise}
	
	\section{Selected Answers}
	
	\begin{itemize}
		\item The number of injective maps
		\[
		\F_q^2\to \F_q^4
		\]
		is
		\[
		(q^4-1)(q^4-q).
		\]
		
		\item Therefore
		\[
		\qbinom{4}{2}
		=
		\frac{(q^4-1)(q^4-q)}{(q^2-1)(q^2-q)}.
		\]
		
		\item For \(q=2\),
		\[
		\qbinom{4}{2}
		=
		\frac{(2^4-1)(2^4-2)}{(2^2-1)(2^2-2)}
		=
		\frac{15\cdot 14}{3\cdot 2}
		=
		35.
		\]
	\end{itemize}
	
	\section*{One-Sentence Takeaway}
	
	Over a finite field \(\F_q\), the \(k\)-dimensional subspaces of an \(n\)-dimensional vector space can be counted by counting injective or surjective linear operators and dividing by the number of automorphisms of the \(k\)-dimensional model space, which yields the Gaussian binomial coefficient \(\qbinom{n}{k}\).
	
\end{document}